% EDITING: type inside the final {...} of each field. Change the shown font size
% (for example, font=10pt) in that same field block when needed.

\GrantSection{9}{Budget, Milestones, and Sustainability}

\GrantWarning{RFA-RM-27-001 instructs applicants to enter only the five-year total direct costs and states that a detailed budget and other budget forms are not required and will not be accepted. Use the planning fields below for funder outreach or other sponsors, but omit any module prohibited by the target solicitation.}

\begin{tabularx}{\linewidth}{X X X}
\textbf{Requested direct costs / year} & \textbf{Total direct costs} & \textbf{Cost sharing}\\[-1pt]
\GrantTableField[id=budget-annual,font=9pt,width=0.95\linewidth,maxchars=60]{\DefaultDirectCostsPerYear} &
\GrantTableField[id=budget-total,font=9pt,width=0.95\linewidth,maxchars=60]{\DefaultTotalDirectCosts} &
\GrantTableField[id=budget-cost-share,font=9pt,width=0.95\linewidth,maxchars=60]{\$0}\\[6pt]
\textbf{Indirect cost treatment} & \textbf{Program income} & \textbf{Currency}\\[-1pt]
\GrantTableField[id=budget-indirect,font=9pt,width=0.95\linewidth,maxchars=100]{Applicable costs by funder} &
\GrantTableField[id=budget-program-income,font=9pt,width=0.95\linewidth,maxchars=60]{\$0} &
\GrantTableField[id=budget-currency,font=9pt,width=0.95\linewidth,maxchars=20]{USD}
\end{tabularx}

\GrantTextArea[
  id=budget-rationale, font=10pt, height=1.3in, maxchars=1650
]{High-level use of funds and cost rationale}{Funds will primarily support protected PI effort; regulatory and clinical-operations personnel; biostatistics and data management; EDC, secure storage, and cybersecurity; investigational pharmacy and specimen handling; site startup, training, and simulation drills; monitoring and independent medical oversight; ^^Jpatient screening and follow-up support; and repository curation, archiving, and dissemination. This cost structure is appropriate because the project's central challenge is safe, verifiable translation of a complex early-phase oncology platform, not basic capital acquisition alone \citep{nih_rfa_rm_27_001,onpremwhippl,phase1guidance}.
}
\GrantTextArea[
  id=milestones-success, font=10pt, height=2.2in, maxchars=2000
]{Milestones, quantitative success criteria, decision gates, and go/no-go thresholds}{Milestone 1: complete all required registrations, FDA strategy finalization, IRB approval, PRS setup, SOP freeze, and successful activation drills before first enrollment. Milestone 2: enroll the first 3 evaluable participants with no treatment-related death, no uncontained cyber-physical safety failure, and acceptable cohort-review findings. Milestone 3: reach 6 and then 9 evaluable participants with cumulative composite primary-endpoint frequency below the prespecified unacceptable boundary and with complete ^^Jadjudicable safety data on at least 95 percent of required primary-window observations. Milestone 4: complete 13 participants, database lock, and integrated Phase 1 review. Go/no-go thresholds: stop or redesign if a treatment-related death occurs; if repeated major robotic or verification-gate failures cannot be corrected; if the composite primary-endpoint rate exceeds the prespecified unacceptable boundary; or ^^Jif recruitment remains below 50 percent of plan for two consecutive quarters despite corrective action. Advance to next-stage planning if safety is acceptable, feasibility is demonstrated, and data ^^Jcompleteness and operational integrity meet the predefined release criteria.
}
\GrantTextArea[
  id=milestones-timeline, font=10pt, height=1.45in, maxchars=1850
]{Five-year timeline and major deliverables}{Year 1: finalize registrations, FDA and IRB submissions, site qualification, PRS setup, SOPs, training, and first-patient activation. Year 2: complete active accrual and all interim safety/cohort reviews. Year 3: finish follow-up, database lock, integrated Phase 1 analysis, publications, and repository releases of protocol-adjacent artifacts and code where shareable. Year 4: refine the platform based on Phase 1 findings, update simulations and verification frameworks, and prepare the randomized comparative next-stage package. Year 5: disseminate results, complete final controlled-access data releases,^^Jmaintain records and repositories, and transition mature components of the platform to Phase 2 or follow-on funding mechanisms \citep{daraxwhipple,phase1guidance,natlpaiplatf,paiautospons}.
}
\GrantTextArea[
  id=milestones-risks, font=10pt, height=2.1in, maxchars=1750
]{Risk register: scientific, enrollment, regulatory, operational, data, supply, and financial risks}{Scientific risk: unacceptable toxicity or infeasible integration of drug and robotic procedure; mitigation: staged safety review and conservative stopping rules. Enrollment risk: slow identification of eligible KRAS G12-mutated resectable PDAC; mitigation: molecular prescreening, tumor-board referral, and backup referral pathways. Regulatory risk: delay in IND/IDE/IRB clearance; mitigation: early submission sequencing and tracked deficiency responses. Operational risk: credentialing gaps, OR scheduling conflict, or device downtime; mitigation: activation checklist, backup scheduling, spare components, and dry runs. Data risk: incomplete source capture, privacy breach, or unstable software versions; mitigation: EDC controls, audit trails, RBAC, backups, and CAPA. Supply risk: drug, assay, or hardware availability constraints; mitigation: qualification of vendors, inventory thresholds, and contingency contracts. ^^JFinancial risk: underestimation of startup, monitoring, or secure-data costs; mitigation: reserve margins, phased spending, and milestone-based reforecasting.
}
\GrantTextArea[
  id=sustainability, font=10pt, height=2.1in, maxchars=1750
]{Sustainability, dissemination, scale-up, follow-on funding, and post-award transition}{Sustainability will come from converting the award's one-time platform build into reusable regulated assets: validated SOPs, repository-based workflows, analysis code, data dictionaries, safety-governance templates, and a decision-ready Phase 2 package \citep{phase1guidance,oncotrialllm,daraxwhipple}. Dissemination will include peer-reviewed manuscripts, ClinicalTrials.gov reporting, controlled-access data release, public protocol-adjacent materials where appropriate, and conference presentations targeted to oncology, surgery, regulatory science, and biomedical AI communities. If Phase 1 milestones are met, scale-up will focus on a comparative randomized study, broader site qualification, and possibly platform adaptation to additional pancreatic or other solid-tumor settings. Follow-on support could include later-stage NIH mechanisms, NCI-aligned cooperative structures, philanthropy, or mission-aligned non-dilutive translational funding. Post-award transition planning will preserve records, repositories, and governance roles so that no critical knowledge is lost between mechanisms.
}
